
\appendix
\section*{Appendix E (Strict, No Fit): Deriving the Geometric Mass Scale $m_0$ from the UBT Action}
\addcontentsline{toc}{section}{Appendix E (Strict, No Fit): Deriving the Geometric Mass Scale $m_0$ from the UBT Action}

\subsection*{E.1 Objective}
In the strict (no-fit) formulation, the absolute mass scale $m_0$ must follow from the \emph{field dynamics} of UBT rather than any measured input. This appendix derives $m_0$ from the stationary value of the UBT action on the $(t,\psi)$ torus for the lowest toroidal harmonic. No anchoring by $m_e$ or other empirical constants is used.

\subsection*{E.2 Setup: UBT Action on the Complex-Time Torus}
Let $\Theta(q,\tau)$ be the biquaternionic field on $\mathbb{C}^5$ with complex time $\tau=t+i\psi$ and biquaternionic metric $\mathcal{G}_{AB}$ adapted to the product torus $S^1_t\times S^1_\psi$ in the $(t,\psi)$-plane (with $g_{AB}:=\Re(\mathcal{G}_{AB})$ the Hermitian projection used for real observables). Consider the UBT action
\begin{equation}
  S[\Theta,\mathcal{G}] \;=\; \int_{\mathcal{M}} \!\mathrm{d}^5x\,\sqrt{|\Re\mathcal{G}|}\;\Big(
     \mathcal{L}_{\text{kin}}[\Theta,\nabla\Theta;\mathcal{G}]
     \;+\; \mathcal{L}_{\text{int}}[\Theta;\mathcal{G}]
     \;+\; \mathcal{L}_{\text{geom}}[\mathcal{G}]\Big)\,,
  \label{eq:UBT-action}
\end{equation}
where $\mathcal{L}_{\text{geom}}$ captures the toroidal curvature contributions that also generate the running of $\alpha(\mu)$ (Appendix~C).
On the $(t,\psi)$ torus we use the real projection $g_{AB}=\Re(\mathcal{G}_{AB})$ with reduced metric
\begin{equation}
  \mathrm{d}s^2 \,=\, R_t^2\,\mathrm{d}\theta_t^2 \,+\, R_\psi^2\,\mathrm{d}\theta_\psi^2
  \qquad (\theta_t,\theta_\psi\in[0,2\pi)),\quad
  \alpha \,=\, \frac{R_t}{R_\psi}\,.
  \label{eq:torus-metric}
\end{equation}

\subsection*{E.3 Mode Ansatz and Energy Density}
For the lowest toroidal harmonic (electron sheet), adopt the separable ansatz
\begin{equation}
 \Theta(q,\tau) \,=\, \Phi(q)\,\exp\!\big(i n\,\theta_\psi \,-\, i \omega\,\theta_t\big)\,,
 \qquad n\in\mathbb{N}\ \text{(baseline $n=1$)}\,.
 \label{eq:mode-ansatz}
\end{equation}
With~\eqref{eq:torus-metric} and $g^{AB}:=\Re(\mathcal{G}^{AB})$ the inverse Hermitian projection, the kinetic term yields the mode energy density
\begin{equation}
 \mathcal{E}_{\text{kin}} \;\propto\;
  g^{\psi\psi}\,(\partial_{\theta_\psi}\Theta)(\partial_{\theta_\psi}\Theta)^\dagger
 + g^{tt}\,(\partial_{\theta_t}\Theta)(\partial_{\theta_t}\Theta)^\dagger
 \;\propto\;
   \frac{n^2}{R_\psi^2} \,+\, \frac{\omega^2}{R_t^2}\,.
 \label{eq:kin-density}
\end{equation}
The interaction and geometric parts contribute curvature-weighted invariants which, to leading order in the two-circle geometry, introduce the same $2\pi$ normalizations that fix the $\beta$-function coefficients in Appendix~C.

\subsection*{E.4 Stationary Energy per Period and Definition of $m_0$}
Define the energy per torus period (one cell) as
\begin{equation}
  \mathcal{E}_\text{cell}(n,\omega;R_t,R_\psi)\;=\;
  A\,\Big(\frac{n^2}{R_\psi^2} + \frac{\omega^2}{R_t^2}\Big)
  \;+\; \mathcal{U}_{\text{geom}}(R_t,R_\psi)\,,
  \label{eq:cell-energy}
\end{equation}
where $A$ is a normalization from $\mathcal{L}_{\text{kin}}$
\textbf{[DERIVED: $A=\hbar^2/(2m_\text{field}R_\psi^2)$, see Appendix~H]}
and $\mathcal{U}_{\text{geom}}$ is the curvature functional induced by $\mathcal{L}_{\text{geom}}$
\textbf{[SKETCH: explicit form from biquaternionic Ricci scalar, see
\texttt{canonical/geometry/biquaternion\_curvature.tex}]}.

\medskip
\noindent\textbf{Reformulated stationary conditions (no circular $\alpha$ input).}
The stationary torus is found by solving the two conditions
\begin{equation}
  \frac{\partial \mathcal{E}_\text{cell}}{\partial R_t} \,=\, 0\,,\qquad
  \frac{\partial \mathcal{E}_\text{cell}}{\partial R_\psi} \,=\, 0\,,
  \label{eq:stationary-conditions}
\end{equation}
\emph{independently}, without imposing $\alpha=R_t/R_\psi$ as a third constraint. The ratio
\begin{equation}
  \alpha_\text{predicted} \;=\; \frac{R_t^\star}{R_\psi^\star}
  \label{eq:alpha-output}
\end{equation}
is then an \textbf{output} of the variational problem, not an input. [ODVOZENO — removes circular dependence of earlier formulation]

\medskip
\noindent\textbf{Explicit solution for the simplest curvature potential.}
For the candidate geometric potential $\mathcal{U}_{\text{geom}}=-C/(R_t R_\psi)$ with $C>0$ [SKETCH — requires derivation of $C$ from $\mathcal{L}_\text{geom}$], the two conditions~\eqref{eq:stationary-conditions} give:
\begin{align}
  \frac{\partial\mathcal{E}_\text{cell}}{\partial R_t}=0 &\;\Rightarrow\;
    \frac{R_t}{R_\psi} = \frac{2A\,\omega^2}{C}\,, \label{eq:stat-Rt}\\
  \frac{\partial\mathcal{E}_\text{cell}}{\partial R_\psi}=0 &\;\Rightarrow\;
    \frac{R_\psi}{R_t} = \frac{2A\,n^2}{C}\,. \label{eq:stat-Rpsi}
\end{align}
Multiplying~\eqref{eq:stat-Rt} and~\eqref{eq:stat-Rpsi} yields the on-shell frequency constraint [ODVOZENO]:
\begin{equation}
  \omega\cdot n \;=\; \frac{C}{2A}\,.
\end{equation}
Substituting the on-shell condition $\omega_\star = m_0 R_t^\star$ (rest-mass dispersion on $S^1_t$) and evaluating at $n=1$ gives:
\begin{equation}
  \alpha_\text{predicted} = \frac{R_t^\star}{R_\psi^\star} = \frac{C}{2A}\,.
  \label{eq:alpha-from-geometry}
\end{equation}
For $\alpha_\text{predicted}=\alpha_\text{exp}\approx1/137.036$ this requires the curvature-to-kinetic ratio $C/A = 2\alpha_\text{exp}\approx0.01460$ [SKETCH — to be derived from $\mathcal{L}_\text{geom}$; see \texttt{tools/m0\_from\_torus.py}].

The cell energy at the stationary point for $n=1$ evaluates to $\mathcal{E}_\text{cell}^\star = A\,m_0^2\,(n^2-1)=0$, yielding a trivial $m_0=0$ from the energy definition alone. A non-trivial mass scale requires either $n>1$ as the reference harmonic or an independent dimensionful input from $\mathcal{L}_\text{geom}$ [SLEPÁ ULIČKA: $n=1$ energy vanishes; documented in \texttt{tools/m0\_from\_torus.py}].

Evaluated at the lowest sheet ($n=1$) and at the self-consistent frequency $\omega=\omega_\star$, the stationary energy defines the geometric mass scale
\begin{equation}
   m_0 \;\equiv\; \frac{1}{2\pi}\,\mathcal{E}_\text{cell}(n{=}1,\omega_\star;R_t^\star,R_\psi^\star)\,,
   \label{eq:m0-definition}
\end{equation}
where $R_t^\star,R_\psi^\star$ solve~\eqref{eq:stationary-conditions}. The $2\pi$ factor normalizes one period on $S^1_t$. No measured masses enter; $m_0$ is fixed by geometry and dynamics alone, once $\mathcal{U}_{\text{geom}}$ is derived from $\mathcal{L}_{\text{geom}}$.

\subsection*{E.5 Electron Sheet and the Fixed-Point}
For the electron ($n=1$), the lepton fixed point (Appendix~D) reads
\begin{equation}
  m_e \,=\, \frac{m_0\,f(1)}{\alpha(m_e)} \;=\; \frac{m_0}{\alpha(m_e)}\,.
  \label{eq:electron-fp}
\end{equation}
In the \emph{strict} setting, $m_0$ from~\eqref{eq:m0-definition} and $\alpha(\mu)$ from Appendix~C close the system without any empirical anchor. For higher leptons ($n=3,9$ baseline), the same $m_0$ produces
\begin{equation}
  m_\mu \,=\, \frac{m_0\,f(3)}{\alpha(m_\mu)}\,,\qquad
  m_\tau \,=\, \frac{m_0\,f(9)}{\alpha(m_\tau)}\,,
\end{equation}
with $f(n)=n^2$ at lowest harmonic.

\subsection*{E.6 Remarks on Uniqueness and Higher Corrections}
The uniqueness of $m_0$ follows from the strict minimization~\eqref{eq:stationary-conditions} once $\mathcal{U}_{\text{geom}}$ is specified by the same curvature principles that yielded $\beta_1$ and $\beta_2$. Higher-harmonic corrections shift $R_t^\star,R_\psi^\star$ and thus $m_0$ at subleading order, preserving the no-fit nature of the construction.

\subsection*{E.7 Summary}
We have derived $m_0$ as the stationary energy per period of the UBT field on the $(t,\psi)$ torus, using only geometry and the UBT action. This closes the strict spectrum without empirical fits: $m_0$ is determined theoretically, $\alpha(\mu)$ is geometric (Appendix~C), and lepton masses follow from the fixed-point equations (Appendix~D).

\subsection*{E.8 Self-Consistency Equation for $\alpha$ via QED Mass Correction}
\label{sec:E8-alpha-selfconsistency}

\noindent\textbf{Status: [HYPOTÉZA — žádný standardní UBT cutoff nepotvrdil $<1\%$ přesnost bez $m_e$ jako vstupu]}

\medskip
\noindent\textbf{Motivation.}
Numerically, $\alpha \approx m_0/(n^\star m_e)$ where $m_0=0.509856\,\text{MeV}$ [CALIBRATED],
$m_e=0.510999\,\text{MeV}$ (PDG), $n^\star=137$ [TOPOLOGICAL INPUT]. The gap
\begin{equation}
  \Delta m = m_e - m_0 = 1.143\,\text{keV}\,,\quad
  \frac{\Delta m}{m_e} = 0.002237 \approx \frac{\alpha}{2\pi}=0.001161
  \quad\text{[not equal]}
\end{equation}
suggests interpreting $\Delta m$ as a QED radiative correction to the geometric mass $m_0$.

\medskip
\noindent\textbf{QED one-loop mass correction.}
In the QED limit ($\partial_\psi=0$, Appendix~D), the one-loop fermion mass correction reads [ODVOZENO from standard QED; recovered in UBT limit]:
\begin{equation}
  \delta m_{\text{QED}} = \frac{3\alpha}{4\pi}\,m_e\,\ln\!\frac{\Lambda^2}{m_e^2}
  = \frac{3\alpha}{2\pi}\,m_e\,\ln\!\frac{\Lambda}{m_e}\,,
  \label{eq:QED-mass-correction}
\end{equation}
where $\Lambda$ is the UV cutoff.  Setting $m_0 = m_e - \delta m_{\text{QED}}$ and $\alpha = m_0/(n^\star m_e)$ gives the self-consistency equation:
\begin{equation}
  \boxed{n^\star\,\alpha \;+\; g(\alpha) \;=\; 1\,,
  \qquad g(\alpha) = \frac{3\alpha}{2\pi}\,\ln\!\frac{\Lambda}{m_e}}
  \label{eq:alpha-selfconsistency}
\end{equation}
with \emph{zero free parameters}, given $n^\star=137$ and a choice of $\Lambda$.

\medskip
\noindent\textbf{Natural UBT cutoffs and numerical results.}
The self-consistency equation~\eqref{eq:alpha-selfconsistency} is solved numerically in \texttt{tools/alpha\_selfconsistency.py}. Results for natural UBT cutoffs:

\begin{center}
\begin{tabular}{|l|l|c|c|}
\hline
\textbf{Cutoff} $\Lambda$ & \textbf{Motivation} & $\alpha^{-1}_\text{sol}$ & \textbf{Error} \\
\hline
$n^\star m_e = 137\,m_e$ & KK tower at winding $n^\star$ [HYPOTHESIS] & 139.35 & 1.66\% \\
$m_e/\alpha$             & classical electron radius [STANDARD]        & 139.36 & 1.67\% \\
$m_e/\sqrt{\alpha}$      & geometric mean scale [HYPOTHESIS]           & 138.18 & 0.83\% \\
$2\,m_e$                 & pair-creation threshold [STANDARD QED]      & 137.33 & 0.22\% \\
\hline
\end{tabular}
\end{center}

\noindent\textbf{Required cutoff.}
Inverting~\eqref{eq:QED-mass-correction} for $\delta m = \Delta m = 1.143\,\text{keV}$ gives [ODVOZENO]:
\begin{equation}
  \ln\!\frac{\Lambda}{m_e} = \frac{\Delta m/m_e}{3\alpha/(2\pi)} = 0.642
  \quad\Rightarrow\quad \Lambda = 1.90\,m_e \approx 971\,\text{keV}\,.
\end{equation}
This lies between $m_e$ and $2m_e$ (pair threshold) and is not a recognised \emph{independent} UBT geometric scale.

\medskip
\noindent\textbf{Assessment.}
The pair-threshold cutoff $\Lambda=2m_e$ gives error $0.22\%$, matching the $0.22\%$ systematic in $m_0$, but it uses $m_e$ directly as input — the independence requirement is not satisfied.  All purely geometric UBT cutoffs produce errors $>0.8\%$, above the target of $<1\%$.

\emph{Conclusion:} Equation~\eqref{eq:alpha-selfconsistency} represents a genuine \emph{self-consistency} structure (not a tautology for any fixed cutoff), but no natural UBT scale has been identified that satisfies it to $<1\%$ without using $m_e$ as input.

\begin{center}
\fbox{\parbox{0.88\linewidth}{\textbf{Classification: [HYPOTÉZA]}\\
$n^\star\alpha + g(\alpha) = 1$ with $g$ from QED in UBT limit.\\
\textbf{What would confirm it:} derivation of $\Lambda \approx 1.90\,m_e$ (or equivalently $\Lambda = 971\,\text{keV}$) from UBT geometry without reference to $m_e$ or $\alpha$.}}
\end{center}
